\item \textbf{¿Por qué no es posible la siguiente situación?}
Una cuña delgada de material transparente ($n = 1.50$) se encuentra colocada en el aire. La cuña tiene una longitud de $0.500\text{ m}$ y una altura máxima de $1.00\text{ mm}$ en su extremo grueso. Al iluminar normalmente con luz monocromática de $632.8\text{ nm}$ y observar desde arriba, se genera un patrón de interferencia en toda su longitud. Un estudiante se dispone a contar las franjas brillantes a lo largo de la cuña, pero se distrae tras contar con total precisión $5000$ franjas brillantes antes de terminar.